%% kapa: exam=1 round=30 year=2019 period=2 form=A subject=appraisal_law subject_ko=감정평가관계법규 seq=15 number=15
%% source: Q-Net 감정평가사 기출문제 (cst00309) artlSeq=5205205 file=extracted/1008/5205205/2019년 제30회 감정평가사 1차시험 문제지/2019년 제30회 감정평가사 1차시험 문제지/2019년도 제30회 감정평가사 1차 2교시 A형.hwp
%% conversion: kordoc 4.24.51 -> tools/gen1.py (auto); review=proofread-sonnet; flags=none
\begin{question}{15}
감정평가 및 감정평가사에 관한 법령상 감정평가사에 관한 설명으로 옳지 않은 것은?
\begin{choices}
  \item 감정평가사 자격이 취소된 후 3년이 경과되지 아니한 사람은 감정평가사가 될 수 없다.
  \item 감정평가사 결격사유는 감정평가사 등록 및 갱신등록의 거부사유가 된다.
  \item 감정평가사는 둘 이상의 감정평가법인 또는 감정평가사사무소에 소속될 수 없다.
  \item 감정평가사 자격이 있는 사람이 국토교통부장관에게 등록하기 위해서는 1년 이상의 실무수습을 마쳐야 한다.
  \item 감정평가사시험에 합격한 사람은 별도의 연수과정을 마치지 않더라도 감정평가사의 자격이 있다.
\end{choices}
\end{question}
\answer{4}
